\chapter{线性系统求解、基本迭代与平滑性质}
\label{ch:iterative-smoothing}

\chapref{ch:discretization}和\chapref{ch:amr-discretization}已经把连续 PDE 变成离散线性系统
\begin{equation}
A\bm u=\bm b.
\end{equation}
从这里开始，问题发生了性质上的变化：离散方程已经给定，接下来要回答的是\textbf{怎样求这个线性系统，以及怎样判断一个求解过程是否真的有效}。在进入 Jacobi、Gauss--Seidel 或 multigrid 之前，本章先建立线性求解方法的基本地图，再深入最简单的平稳迭代，并用误差传播、谱半径、收敛因子和网格依赖性解释它们为什么既“能收敛”，又可能随着网格加密迅速变慢。

本章仍然以一维 Poisson 为最小可检查模型，但分析对象已经提升到代数层。除非特别指出结论依赖局部守恒、对称性、stencil 几何或系数位置，否则误差、残差、stationary iteration 和谱收敛并不要求 $A$ 必须来自 FD 或 FV 中某一种特定构造。

\begin{keybox}
\textbf{本章要建立的主线：}
\[
\boxed{
A\bm u=\bm b
\rightarrow
\text{求解器类型}
\rightarrow
\text{误差/残差}
\rightarrow
\text{局部迭代}
\rightarrow
\text{误差传播}
\rightarrow
\text{收敛速度}
\rightarrow
\text{网格依赖退化}
\rightarrow
\text{平滑性质}
}
\]
第四章的 coarse-grid correction 将由本章最后暴露出的“低频误差难以被局部迭代消除”自然触发。
\end{keybox}

\section{离散线性系统的求解方法地图}
\label{sec:linear-solver-map}

面对
\begin{equation}
A\bm u=\bm b,
\end{equation}
“求解矩阵”并不只有 Jacobi 或 Gauss--Seidel 一条路线。对 PDE 离散产生的大型稀疏系统，至少应先区分四类基本策略；预条件则是横跨多类迭代方法的加速机制。

\subsection{直接法}

直接法通过有限次代数变换把系统化为更容易求解的形式。对一般非奇异矩阵，实际高斯消元通常需要行置换，可写成
\begin{equation}
PA=LU,
\end{equation}
其中 $P$ 是置换矩阵；随后依次求解下三角和上三角系统。若 $A$ 对称正定，则不需要这种行主元选取，可使用 Cholesky 分解
\begin{equation}
A=LL^T.
\end{equation}

“直接”并不意味着浮点计算中绝对没有误差，而是指算法不依赖“第 $k$ 轮近似逐渐趋近真解”这一收敛过程。在精确算术下，分解与回代经过有限步骤即可得到解。

对 PDE 稀疏矩阵，直接法的主要困难是\textbf{fill-in}：原矩阵中为零的位置在消元后可能产生新的非零元，使因子 $L/U$ 比原始 stencil 矩阵稠密得多。因此直接法非常适合中小规模问题、粗层问题或需要重复求解且因子可以复用的场景，但在大规模三维 PDE 上，内存和分解成本可能成为限制。

\subsection{平稳迭代}

另一类方法不一次完成消元，而是从初始猜测 $\bm u^{(0)}$ 出发，重复固定的更新规则。很多方法都可以写成
\begin{equation}
\bm u^{(k+1)}
=
\bm u^{(k)}
+
M^{-1}\left(\bm b-A\bm u^{(k)}\right),
\label{eq:iteration-stationary-preview}
\end{equation}
其中 $M$ 在每一轮保持不变，并且比完整的 $A$ 更容易处理。

Jacobi、weighted Jacobi、Gauss--Seidel 和 SOR 都属于\textbf{平稳迭代（stationary iteration）}。本章会重点研究这一类，因为 Jacobi/GS 不仅可以独立求解小问题，更重要的是它们会成为 multigrid 中的 smoother。

\subsection{Krylov 子空间方法}

Krylov 方法也从初始残差
\begin{equation}
\bm r_0=\bm b-A\bm u^{(0)}
\end{equation}
出发，但它不会每轮机械重复同一个固定修正，而是逐步利用
\begin{equation}
\mathcal K_m(A,\bm r_0)
=
\operatorname{span}
\left\{
\bm r_0,\,
A\bm r_0,\,
A^2\bm r_0,\ldots,
A^{m-1}\bm r_0
\right\}
\label{eq:iteration-krylov-space}
\end{equation}
中累积的信息构造新的近似。

此处只需要建立最低可用认识：Conjugate Gradient（CG）是 SPD 系统的重要代表；GMRES 面向更一般的非对称线性系统。它们的搜索方向和历史信息组织比 stationary iteration 更丰富。后续\chapref{ch:mg-cycles}会在 multigrid 已经形成近似逆以后，正式讨论 CG/GMRES/FGMRES 与 multigrid preconditioner 的组合。

\subsection{多重网格与预条件}

Multigrid 与前两类迭代的核心差异，是显式利用误差的\textbf{尺度结构}：局部平滑器处理细网格上变化快的误差，coarse-grid correction 处理在细网格上变化慢的误差。本章只先建立前半部分的 smoother 机制；\chapref{ch:two-grid}和\chapref{ch:mg-cycles}分别完成 coarse-grid correction 与多层递归。

预条件（preconditioning）则不应被理解成与上述方法平行的“第五种求解器”。它的目标是构造一个容易应用的近似逆
\begin{equation}
M^{-1}\approx A^{-1},
\end{equation}
把原问题变成对迭代算法更友好的形式。Jacobi 的 $D^{-1}$ 可以看成极简单的近似逆；一个完整 V-cycle 则可以成为更强的 multigrid preconditioner。

\begin{table}[htbp]
\centering
\small
\begin{tabularx}{0.96\textwidth}{>{\raggedright\arraybackslash}p{3.0cm}
                                  >{\raggedright\arraybackslash}p{4.5cm}
                                  >{\centering\arraybackslash}p{2.2cm}
                                  >{\raggedright\arraybackslash}X}
\toprule
方法类别 & 核心动作 & 主要信息结构 & 本书中的主要角色\\
\midrule
LU / Cholesky & 分解后做三角求解 & 稀疏因子 $L/U$ & 小规模问题、bottom solve、基准对照\\
Jacobi / GS & 重复固定局部修正规则 & 固定近似矩阵 $M$ & 基本 solver、multigrid smoother\\
CG / GMRES & 在 Krylov 子空间中构造搜索方向 & Krylov 历史与内积 & 大规模外层迭代求解器\\
Multigrid & 在不同尺度上处理不同误差成分 & 层级、transfer 与 coarse operator & solver 或 preconditioner\\
\bottomrule
\end{tabularx}
\caption{本书涉及的主要线性求解路线。这里给出的是角色地图，不替代后续各算法的正式推导。}
\label{tab:iteration-solver-map}
\end{table}

\begin{keybox}
\textbf{本章的范围边界：}
这里不展开 LU pivoting、完整 CG recurrence 或 Arnoldi/GMRES 推导。本章深入 stationary iteration，是因为它既能建立最基本的线性收敛理论，又直接解释 multigrid smoother 为什么有效、为什么又不够。
\end{keybox}

\section{误差、残差与修正方程}

设真解为 $\bm u^\star$，当前近似为 $\bm u^{(k)}$。定义
\begin{equation}
\bm e^{(k)}
=
\bm u^\star-\bm u^{(k)}
\label{eq:iteration-error-def}
\end{equation}
为\textbf{误差（error）}。它直接表示当前解离真解有多远，但实际计算中通常不知道 $\bm u^\star$，因此无法直接计算 $\bm e^{(k)}$。

把当前近似代回原方程，定义
\begin{equation}
\bm r^{(k)}
=
\bm b-A\bm u^{(k)}
\label{eq:iteration-residual-def}
\end{equation}
为\textbf{残差（residual）}。因为 $A\bm u^\star=\bm b$，
\begin{equation}
\bm r^{(k)}
=
A\left(\bm u^\star-\bm u^{(k)}\right)
=
A\bm e^{(k)}.
\label{eq:iteration-residual-error-relation}
\end{equation}
因此残差不是误差本身，而是误差经过算子 $A$ 作用后的结果。

式~\eqref{eq:iteration-residual-error-relation} 还可以反过来读成\textbf{修正方程}：
\begin{equation}
A\,\delta\bm u
=
\bm r^{(k)}.
\label{eq:iteration-correction-equation}
\end{equation}
若能够精确求解它，就有
\begin{equation}
\delta\bm u
=
A^{-1}\bm r^{(k)}
=
\bm e^{(k)},
\end{equation}
于是
\begin{equation}
\bm u^{(k+1)}
=
\bm u^{(k)}+\delta\bm u
=
\bm u^\star.
\end{equation}
所以许多迭代方法都可以理解成：\textbf{用一个比精确求解 $A\delta\bm u=\bm r$ 更便宜的过程，近似构造误差修正。}

例如
\begin{equation}
A=
\begin{bmatrix}
2&-1&0\\
-1&2&-1\\
0&-1&2
\end{bmatrix},
\qquad
\bm b=
\begin{bmatrix}
1\\1\\1
\end{bmatrix},
\qquad
\bm u^{(0)}=\bm0.
\end{equation}
即使不知道真解，也可以立即得到
\begin{equation}
\bm r^{(0)}=\bm b.
\end{equation}
这就是迭代程序实际能够直接观测的量。

\subsection{小残差是否一定意味着小误差}

若 $A$ 可逆，由
\begin{equation}
\bm e=A^{-1}\bm r
\end{equation}
可得
\begin{equation}
\|\bm e\|
\le
\|A^{-1}\|\,\|\bm r\|.
\label{eq:iteration-error-residual-bound}
\end{equation}
因此残差怎样映射成误差，取决于 $A^{-1}$ 的放大能力。

定义矩阵条件数
\begin{equation}
\kappa(A)
=
\|A\|\,\|A^{-1}\|.
\label{eq:iteration-condition-number}
\end{equation}
当 $\bm b=A\bm u^\star\neq0$ 时，可以得到相对量的基本上界
\begin{equation}
\frac{\|\bm e\|}{\|\bm u^\star\|}
\le
\kappa(A)
\frac{\|\bm r\|}{\|\bm b\|}.
\label{eq:iteration-relative-residual-error}
\end{equation}
这并不是说条件数完全决定每一次实际误差，而是说明一个重要事实：\textbf{同样大小的相对残差，在病态系统中可能对应更大的解误差。}

纯 Neumann 问题存在常数零空间时，$A^{-1}$ 在整个空间上不存在；此时必须先进入满足兼容条件并固定 gauge 的可逆子空间，再谈类似的误差--残差关系。

\begin{keybox}
\textbf{三个对象要始终分开：}
\[
\boxed{
\text{error：离真解多远}
\qquad
\text{residual：方程满足得多好}
\qquad
\text{conditioning：二者之间可能放大多少}
}
\]
\end{keybox}

\section{Jacobi 迭代}

\subsection{Jacobi 点迭代公式}

回到一维 Poisson 离散，把共同的 $1/h^2$ 因子乘到右端，一个内部点满足
\begin{equation}
2u_i-u_{i-1}-u_{i+1}
=
h^2f_i.
\label{eq:iteration-point-equation}
\end{equation}
把中心未知量单独写出：
\begin{equation}
u_i
=
\frac12
\left(
h^2f_i+u_{i-1}+u_{i+1}
\right).
\label{eq:jacobi_point}
\end{equation}

如果右侧邻居还不是真解，就用上一轮近似代替。Jacobi 定义为
\begin{equation}
u_i^{(k+1)}
=
\frac12
\left(
h^2f_i
+
u_{i-1}^{(k)}
+
u_{i+1}^{(k)}
\right).
\label{eq:iteration-jacobi-one-dimensional}
\end{equation}
同一轮中所有点只读旧状态 $\bm u^{(k)}$，把结果写入新状态 $\bm u^{(k+1)}$。因此一轮内部没有逐点先后依赖。

\subsection{三未知量离散系统}

取三个内部未知量，两端 Dirichlet 边界为零，并令
\begin{equation}
h^2f_i=1,
\qquad i=1,2,3.
\end{equation}
从
\begin{equation}
\bm u^{(0)}=(0,0,0)^T
\end{equation}
开始，第一轮得到
\begin{equation}
\bm u^{(1)}
=
\left(
\frac12,\frac12,\frac12
\right)^T,
\end{equation}
第二轮得到
\begin{equation}
\bm u^{(2)}
=
\left(
\frac34,1,\frac34
\right)^T.
\end{equation}

这展示了迭代求解最基本的机制：不是一次直接构造 $\bm u^\star$，而是重复一个便宜的局部规则，让近似逐步变化。

\begin{checkpointbox}
\textbf{即时检查 3.1}

继续计算 $\bm u^{(3)}$。这一轮三个新值为什么必须全部读取 $\bm u^{(2)}$，而不能读取本轮刚算出的其他分量？
\end{checkpointbox}

\section{Gauss--Seidel 迭代}

Jacobi 明确区分旧数组和新数组。另一种自然选择是：如果某个邻居在本轮已经算出新值，就立即使用它。

对一维 Poisson，若按 $i=1,2,\ldots,N$ 从左到右扫描，
\begin{equation}
u_i^{(k+1)}
=
\frac12
\left(
h^2f_i
+
u_{i-1}^{(k+1)}
+
u_{i+1}^{(k)}
\right).
\label{eq:iteration-gs-one-dimensional}
\end{equation}
左邻居已经属于本轮新状态，右邻居仍使用旧状态。

对上一节三未知量例子，从零向量开始做一次 sweep：
\begin{align}
u_1^{(1)}&=\frac12,\\
u_2^{(1)}&=\frac12\left(1+u_1^{(1)}+u_3^{(0)}\right)=\frac34,\\
u_3^{(1)}&=\frac12\left(1+u_2^{(1)}\right)=\frac78.
\end{align}
所以
\begin{equation}
\bm u^{(1)}_{GS}
=
\left(
\frac12,\frac34,\frac78
\right)^T.
\end{equation}

这里最重要的不是第一轮“谁更接近真解”，而是数据依赖已经改变：Jacobi 一轮中的点更新彼此独立；lexicographic Gauss--Seidel 则具有明确扫描顺序。

\subsection{红黑 Gauss--Seidel}

二维五点网格可以按 $(i+j)\bmod2$ 染成红、黑两色。同一颜色的点之间没有直接 stencil 耦合，因此可以先更新全部红点，再更新全部黑点。这就是 red--black Gauss--Seidel（RBGS）。

\begin{algorithm}[htbp]
\caption{Red--Black Gauss--Seidel 单次 sweep}
\begin{algorithmic}[1]
\For{$c\in\{\text{red},\text{black}\}$}
  \State 更新所有颜色为 $c$ 的网格点
  \State 完成颜色 $c$ 后再进入下一颜色
\EndFor
\end{algorithmic}
\end{algorithm}

RBGS 没有消除 Gauss--Seidel 的依赖，而是把逐点依赖重组为两个颜色阶段。\chapref{ch:shared-memory}的\secref{sec:shared-rbgs}会把这种数学依赖进一步解释成线程同步成本。

\section{多维网格上的迭代结构}

对二维 Poisson，Jacobi 更新为
\begin{equation}
u_{i,j}^{(k+1)}
=
\frac{1}{2(h_x^{-2}+h_y^{-2})}
\left[
f_{i,j}
+h_x^{-2}
\left(
u_{i-1,j}^{(k)}+u_{i+1,j}^{(k)}
\right)
+h_y^{-2}
\left(
u_{i,j-1}^{(k)}+u_{i,j+1}^{(k)}
\right)
\right].
\label{eq:iteration-two-dimensional-jacobi}
\end{equation}
三维再增加 $z$ 向两个邻居。Jacobi 的关键性质仍然不变：同一 sweep 的输出全部只读旧状态，因此二维、三维仍然具有 pointwise 并行结构。

Gauss--Seidel 则不同。二维字典序扫描时，$(i-1,j)$ 和通常的 $(i,j-1)$ 已经是新值，而 $(i+1,j)$、$(i,j+1)$ 仍是旧值；进入三维以后，还会出现当前切片之前/之后的依赖。因此 lexicographic GS 的方向偏置会随着维数增加而更加明显。

对标准 Cartesian 五点/七点 stencil，可以按坐标奇偶性做二着色；但对更宽 stencil、非结构图或 AMR coarse--fine coupling，二着色是否仍然足够必须重新检查 matrix graph，不能从五点/七点格式机械外推。

若
\begin{equation}
-\epsilon u_{xx}-u_{yy}=f,
\qquad
\epsilon\ll1,
\end{equation}
不同方向的耦合强度又会严重不均衡，point smoother 可能无法有效处理沿强耦合方向组织的误差。line/plane smoother、semicoarsening 等方法由此产生，定量分析见\chapref{ch:lfa}。

\section{迭代矩阵与误差传播}

现在把具体更新抽象成统一形式。写
\begin{equation}
A=D+L+U,
\end{equation}
其中 $D$ 是对角部分，$L$、$U$ 分别为严格下三角和严格上三角部分。

一般 stationary iteration 为
\begin{equation}
\bm u^{(k+1)}
=
\bm u^{(k)}
+
M^{-1}
\left(
\bm b-A\bm u^{(k)}
\right),
\label{eq:iteration-stationary-general}
\end{equation}
其中 $M$ 固定不变且容易求解。Jacobi 取
\begin{equation}
M=D,
\end{equation}
Gauss--Seidel 取
\begin{equation}
M=D+L.
\end{equation}

用真解减去式~\eqref{eq:iteration-stationary-general} 两边：
\begin{align}
\bm e^{(k+1)}
&=
\bm e^{(k)}
-
M^{-1}A\bm e^{(k)}
\\
&=
\left(
I-M^{-1}A
\right)\bm e^{(k)}.
\end{align}
定义
\begin{equation}
S
=
I-M^{-1}A
\label{eq:error_prop_general}
\end{equation}
为\textbf{误差传播矩阵（error-propagation matrix）}，则
\begin{equation}
\bm e^{(k)}
=
S^k\bm e^{(0)}.
\label{eq:iteration-error-power}
\end{equation}

这一步非常关键：我们不再只问“更新公式长什么样”，而是开始问
\begin{equation}
\boxed{
S^k\bm e^{(0)}
\quad
\text{是否趋于零，以及趋于零有多快？}
}
\end{equation}

\section{收敛性、收敛阶与收敛因子}
\label{sec:iteration-convergence}

\subsection{谱半径决定是否最终收敛}

若
\begin{equation}
S\bm v=\lambda\bm v,
\end{equation}
则这个特征方向经过 $k$ 轮以后变成
\begin{equation}
S^k\bm v
=
\lambda^k\bm v.
\end{equation}
所以 $|\lambda|<1$ 时该分量衰减，$|\lambda|>1$ 时增长，$|\lambda|=1$ 时幅值不会消失。

定义谱半径
\begin{equation}
\rho(S)
=
\max_i|\lambda_i|.
\end{equation}
对于有限维线性 stationary iteration，从任意初始误差出发都收敛到真解的充要条件是
\begin{equation}
\boxed{
\rho(S)<1.
}
\label{eq:iteration-spectral-convergence}
\end{equation}

这里不能把“$\rho(S)<1$”误读成“每一轮误差范数都严格缩小”。若 $S$ 非正规（non-normal），某些范数下可能出现短暂增长；但只要 $\rho(S)<1$，最终仍有 $S^k\to0$。对任意相容矩阵范数，Gelfand 公式给出
\begin{equation}
\lim_{k\to\infty}
\|S^k\|^{1/k}
=
\rho(S),
\label{eq:iteration-gelfand}
\end{equation}
因此谱半径控制的是最坏情形的渐近指数衰减尺度。

\subsection{“收敛阶”与“收敛因子”不是一回事}

一般非线性迭代中，经典收敛阶 $p$ 可以由
\begin{equation}
\lim_{k\to\infty}
\frac{
\|\bm e^{(k+1)}\|
}{
\|\bm e^{(k)}\|^p
}
=
C
\label{eq:iteration-convergence-order}
\end{equation}
定义。$p=1$ 称为线性收敛，$p=2$ 称为二次收敛；当 $p=1$ 时通常还要求 $0<C<1$。

Jacobi、Gauss--Seidel 这类固定线性 stationary iteration 的典型渐近行为本来就是
\begin{equation}
\|\bm e^{(k)}\|
\sim
q^k\|\bm e^{(0)}\|,
\qquad
0<q<1,
\label{eq:iteration-geometric-convergence}
\end{equation}
所以仅仅说“它们是一阶收敛”区分度很低。更有信息量的是\textbf{收敛因子（convergence factor）} $q$ 有多大。

若渐近阶段可以用固定 $q$ 近似，并希望把误差降到初始值的 $\varepsilon$ 倍，则
\begin{equation}
q^k\approx\varepsilon,
\end{equation}
于是
\begin{equation}
\boxed{
k
\approx
\frac{\log\varepsilon}{\log q}.
}
\label{eq:iteration-count-from-factor}
\end{equation}
当 $q$ 接近 1 时，哪怕仍然满足“收敛”，所需迭代数也可能非常大。

还必须把这里的“迭代收敛阶”与 PDE 离散中的网格收敛阶区分开。二阶空间离散
\begin{equation}
\|u-u_h\|=\Oh(h^2)
\end{equation}
描述的是 $h\to0$ 时离散解逼近连续解；式~\eqref{eq:iteration-geometric-convergence} 描述的是固定离散系统上 $k\to\infty$ 时代数误差的衰减。两者的自变量、误差来源和评价方法都不同。

\begin{checkpointbox}
\textbf{即时检查 3.2}

若两个迭代法的渐近收敛因子分别为 $q=0.2$ 和 $q=0.95$，把误差降低到 $10^{-6}$ 倍大约各需要多少轮？为什么“二者都满足 $q<1$”并不能说明它们同样有效？
\end{checkpointbox}

\section{Jacobi 的网格依赖收敛}
\label{sec:jacobi-grid-dependent}

现在对一维 Dirichlet Poisson 做一个完整的谱计算。设有 $N$ 个内部点，
\begin{equation}
A
=
\frac1{h^2}
\begin{bmatrix}
2&-1&&\\
-1&2&-1&\\
&\ddots&\ddots&\ddots\\
&&-1&2
\end{bmatrix},
\end{equation}
则
\begin{equation}
D=\frac2{h^2}I,
\end{equation}
Jacobi 误差传播矩阵为
\begin{equation}
S_J
=
I-D^{-1}A
=
\frac12
\begin{bmatrix}
0&1&&\\
1&0&1&\\
&\ddots&\ddots&\ddots\\
&&1&0
\end{bmatrix}.
\label{eq:iteration-jacobi-error-matrix}
\end{equation}

取离散正弦模态
\begin{equation}
v_i^{(j)}
=
\sin(i\theta_j),
\qquad
\theta_j
=
\frac{j\pi}{N+1},
\qquad
j=1,\ldots,N.
\label{eq:iteration-jacobi-eigenmode}
\end{equation}
由恒等式
\begin{equation}
\frac12
\left[
\sin((i-1)\theta)
+
\sin((i+1)\theta)
\right]
=
\cos\theta\,\sin(i\theta),
\end{equation}
得到
\begin{equation}
S_J\bm v^{(j)}
=
\lambda_j\bm v^{(j)},
\qquad
\lambda_j
=
\cos
\left(
\frac{j\pi}{N+1}
\right).
\label{eq:iteration-jacobi-spectrum}
\end{equation}
所以
\begin{equation}
\rho(S_J)
=
\cos
\left(
\frac{\pi}{N+1}
\right).
\label{eq:iteration-jacobi-rho}
\end{equation}

当 $N$ 很大，
\begin{equation}
\cos x
=
1-\frac{x^2}{2}
+\Oh(x^4),
\end{equation}
于是
\begin{equation}
\rho(S_J)
=
1-
\frac{\pi^2}{2(N+1)^2}
+
\Oh(N^{-4}).
\label{eq:iteration-jacobi-rho-asymptotic}
\end{equation}
在单位区间上 $h=1/(N+1)$，因此
\begin{equation}
\boxed{
1-\rho(S_J)
=
\Theta(h^2).
}
\label{eq:iteration-jacobi-gap-h2}
\end{equation}

这意味着网格越细，
\begin{equation}
\rho(S_J)\to1.
\end{equation}
对固定误差降低目标 $\varepsilon$，由式~\eqref{eq:iteration-count-from-factor} 与 $\log(1-c h^2)\sim-c h^2$，
\begin{equation}
k
=
\Theta
\left(
h^{-2}\log\frac1\varepsilon
\right).
\label{eq:iteration-jacobi-iteration-complexity}
\end{equation}
若 $\varepsilon$ 固定，就是
\begin{equation}
\boxed{k=\Theta(h^{-2}).}
\end{equation}

所以 Jacobi 作为独立 solver 在细网格上变慢，不是因为某个实现“写得不好”，而是因为最慢误差模态的放大因子随着网格加密系统性逼近 1。

对这个标准一维模型，lexicographic Gauss--Seidel 有
\begin{equation}
\rho(S_{GS})
=
\cos^2
\left(
\frac{\pi}{N+1}
\right)
=
1-
\frac{\pi^2}{(N+1)^2}
+
\Oh(N^{-4}),
\label{eq:iteration-gs-rho-model}
\end{equation}
因此它改善了常数，却没有改变
\begin{equation}
k=\Theta(h^{-2})
\end{equation}
这一网格依赖的量级。

\begin{keybox}
\textbf{这是 multigrid 的第一个关键动机：}
\[
\boxed{
\text{局部 stationary iteration 可以收敛，}
\quad
\text{但其全局收敛速度会随网格加密恶化。}
}
\]
接下来要进一步问：到底是哪一类误差拖慢了它？
\end{keybox}

\section{加权 Jacobi}

\subsection{普通 Jacobi 的高频问题}

对齐次误差方程，普通 Jacobi 在一维满足
\begin{equation}
e_i^{(k+1)}
=
\frac12
\left(
e_{i-1}^{(k)}
+
e_{i+1}^{(k)}
\right).
\label{eq:iteration-jacobi-error-point}
\end{equation}
考虑交替误差
\begin{equation}
e_i^{(k)}=(-1)^i.
\label{eq:iteration-alternating-error}
\end{equation}
它的两个邻居都等于 $-e_i^{(k)}$，因此
\begin{equation}
e_i^{(k+1)}
=
-e_i^{(k)}.
\end{equation}
误差只翻转符号，幅值没有减小。

这与上一节的谱完全一致：最高频离散模态对应的 Jacobi 特征值接近 $-1$。所以普通 Jacobi 虽然整体满足 $\rho(S_J)<1$，却不是理想的 multigrid smoother。

\subsection{阻尼参数与加权误差传播}

设普通 Jacobi 候选值为 $u_i^J$。Weighted Jacobi 定义
\begin{equation}
u_i^{(k+1)}
=
(1-\omega)u_i^{(k)}
+
\omega u_i^J,
\qquad
0<\omega\le1.
\label{eq:iteration-weighted-blend}
\end{equation}
矩阵形式为
\begin{equation}
\bm u^{(k+1)}
=
\bm u^{(k)}
+
\omega D^{-1}
\left(
\bm b-A\bm u^{(k)}
\right),
\label{eq:iteration-weighted-jacobi}
\end{equation}
因此误差传播算子为
\begin{equation}
S_\omega
=
I-\omega D^{-1}A.
\label{eq:iteration-weighted-error-matrix}
\end{equation}

对交替误差，普通 Jacobi 候选值恰好为 $-e_i^{(k)}$，所以
\begin{equation}
e_i^{(k+1)}
=
(1-2\omega)e_i^{(k)}.
\label{eq:iteration-weighted-alt}
\end{equation}
若取
\begin{equation}
\omega=\frac23,
\end{equation}
则
\begin{equation}
e_i^{(k+1)}
=
-\frac13e_i^{(k)}.
\end{equation}

更一般地，对上一节的正弦模态，weighted Jacobi 的放大因子为
\begin{equation}
\mu_\omega(\theta)
=
1-\omega(1-\cos\theta).
\label{eq:iteration-weighted-mode-factor}
\end{equation}
若先把相对于 $2h$ 粗网格不可独立解析的高频区粗略记为
\begin{equation}
\frac{\pi}{2}\le|\theta|\le\pi,
\end{equation}
则
\begin{equation}
1-\cos\theta\in[1,2].
\end{equation}
选择 $\omega=2/3$ 时，
\begin{equation}
\max_{\pi/2\le|\theta|\le\pi}
|\mu_\omega(\theta)|
=
\frac13.
\label{eq:iteration-weighted-high-frequency-factor}
\end{equation}
也就是说，它并不是为了让所有频率都以同一个速度收敛，而是专门把高频区的最坏放大因子压下来。\chapref{ch:lfa}会在无限/周期局部模型上把这一分析正式写成 Fourier symbol，并继续构造 two-grid convergence factor。

\begin{checkpointbox}
\textbf{即时检查 3.3}

对交替误差分别计算 $\omega=1$、$\omega=1/2$ 和 $\omega=2/3$ 的放大因子。哪一个能一步消掉这一单独模式？为什么“对一个模式最优”仍不等于“对整个高频区最优”？
\end{checkpointbox}

\section{平滑性质}
\label{sec:smoothing-property}

现在把“全局收敛”和“作为 smoother 是否有效”分开。

一个独立 solver 关心整个误差空间中的最慢衰减方向，因此自然看
\begin{equation}
\rho(S)
\end{equation}
或相应的全局 convergence factor。Multigrid smoother 的任务却更窄：它只需要快速削弱\textbf{细网格上振荡较快的误差}。

对于缓慢变化的误差，
\begin{equation}
e_{i-1}\approx e_i\approx e_{i+1},
\end{equation}
邻居平均本来就与中心值接近，因此局部更新产生的修正很小。对快速振荡误差，中心点与邻居之间差异大，合适的 weighted Jacobi、Gauss--Seidel 或更强局部方法可以在少量 sweeps 中显著削弱其幅值。经过若干轮以后，剩余误差通常在空间上看起来更加光滑，这就是\textbf{平滑性质（smoothing property）}。

因此必须区分两个评价对象：
\begin{equation}
\boxed{
\text{global convergence factor}
\neq
\text{smoothing factor}.
}
\end{equation}
前者问“整个系统作为独立 solver 收敛多快”；后者只问“被 coarse grid 视为高频的那部分误差压得多快”。一个方法完全可能不是好的独立 solver，却是很好的 smoother。

这也解释了为什么上一节并不试图让 weighted Jacobi 把低频模态也迅速压掉。低频误差正是下一章 coarse-grid correction 要接手的对象。

\begin{keybox}
\textbf{从本章到 two-grid 的逻辑闭环：}
\[
\boxed{
\begin{aligned}
&\text{局部迭代可收敛，但 } \rho(S)\to1 \text{ 随网格加密恶化};\\
&\text{合适 smoother 能快速压低高频误差};\\
&\text{剩余低频误差在细网格上仍然衰减很慢};\\
&\Longrightarrow\text{需要一个能直接处理低频误差的粗空间。}
\end{aligned}
}
\]
\end{keybox}

\section{本章小结}

本章第一次把“离散以后怎样求解”完整放在线性代数视角下。直接法通过分解和三角求解完成有限步骤消元；stationary iteration 重复固定修正规则；Krylov 方法利用不断扩展的 $\mathcal K_m(A,r_0)$ 和历史搜索信息；multigrid 则利用误差的尺度结构。预条件是横跨迭代方法的近似逆机制，而不是与这些方法并列的独立类别。

对 stationary iteration，误差与残差满足
\begin{equation}
\bm r=A\bm e,
\end{equation}
而 conditioning 决定“小残差”能够多可靠地转化成“小误差”。一般误差传播写成
\begin{equation}
\bm e^{(k+1)}
=
S\bm e^{(k)},
\end{equation}
有限维线性 stationary iteration 从任意初值收敛的核心条件是
\begin{equation}
\rho(S)<1.
\end{equation}
但“是否收敛”还不够：真正决定迭代次数的是渐近收敛因子。对一维 Poisson Jacobi，
\begin{equation}
1-\rho(S_J)
=
\Theta(h^2),
\end{equation}
所以固定误差降低目标需要的迭代数按 $\Theta(h^{-2})$ 增长。

这个网格依赖退化并不意味着 Jacobi/GS 对所有误差都同样无效。Weighted Jacobi 可以把高频误差快速压低，而留下在空间上更光滑的误差。于是下一章不再要求 smoother 独立解决整个系统，而是让 coarse-grid correction 专门处理这些局部迭代最难消除的低频成分。

\section*{习题}

本章习题围绕四种能力展开：识别求解路线、重构迭代机制、计算收敛因子、诊断网格依赖性。

\subsection*{求解器地图与迭代机制}
\begin{enumerate}[label=\arabic*.]
\item 对以下任务分别说明更接近 direct、stationary、Krylov 还是 multigrid 思路，并解释理由：
  (a) 小型 SPD 粗层系统反复求解；
  (b) 大型 SPD 稀疏系统只提供 operator apply；
  (c) 非对称大型系统配合一个固定 V-cycle；
  (d) 只希望用两次廉价局部 sweep 压低高频误差。
\item 对
\[
A=
\begin{bmatrix}
2&-1&0\\
-1&2&-1\\
0&-1&2
\end{bmatrix},
\qquad
\bm b=
\begin{bmatrix}
1\\0\\1
\end{bmatrix},
\qquad
\bm u^{(0)}=\bm0,
\]
手算三轮 Jacobi，列出每轮 $\bm u^{(k)}$ 与 $\bm r^{(k)}$。
\item 对上一题手算一轮从左到右 Gauss--Seidel 和一轮从右到左 Gauss--Seidel，指出顺序依赖出现在哪里。
\item 从
\[
\bm u^{(k+1)}
=
\bm u^{(k)}+M^{-1}(\bm b-A\bm u^{(k)})
\]
推导 $S=I-M^{-1}A$，并分别写出 Jacobi 与 lexicographic Gauss--Seidel 的 $M$。
\item 已知精确解 $\bm u^\star=(1,1,1)^T$、当前近似 $\bm u=(0.9,1.1,0.8)^T$，计算误差与残差并验证 $\bm r=A\bm e$。
\end{enumerate}

\subsection*{收敛因子、条件数与网格依赖}
\begin{enumerate}[label=\arabic*.]
\item 某两种 stationary iteration 的渐近收敛因子分别为 $0.2$ 与 $0.95$。用式~\eqref{eq:iteration-count-from-factor}估计把误差降低 $10^{-8}$ 所需的迭代数。
\item 从 $\bm e=A^{-1}\bm r$ 推导式~\eqref{eq:iteration-relative-residual-error}。这个上界为什么不能反过来说明“小误差一定要求极小残差”？
\item 对一维 $N$ 点 Dirichlet Poisson，使用正弦模态完整推导式~\eqref{eq:iteration-jacobi-spectrum} 与谱半径式~\eqref{eq:iteration-jacobi-rho}。
\item 由
\[
\rho(S_J)
=
1-\frac{\pi^2}{2(N+1)^2}
+\Oh(N^{-4})
\]
推导固定 reduction factor 下 $k=\Theta(N^2)$，并转换成 $k=\Theta(h^{-2})$。
\item 比较 $N=31,63,127,255$ 时 Jacobi 理论谱半径和把误差降低 $10^{-6}$ 所需的渐近迭代数。观察网格翻倍时迭代数如何增长。
\item 对标准一维模型比较 $\rho(S_J)$ 与 $\rho(S_{GS})$ 的渐近展开。为什么 Gauss--Seidel 虽然更快，却仍没有消除 mesh dependence？
\item 证明在有限维线性空间中 $\rho(S)<1$ 当且仅当 $S^k\to0$；可以使用 Jordan 标准形或你熟悉的等价结论。
\end{enumerate}

\subsection*{平滑性质与数值验证}
\begin{enumerate}[label=\arabic*.]
\item 对 weighted Jacobi 推导
\[
\mu_\omega(\theta)
=
1-\omega(1-\cos\theta).
\]
在 $\theta\in[\pi/2,\pi]$ 上比较 $\omega=1/2,2/3,1$ 的最坏放大因子。
\item 实现 Jacobi、weighted Jacobi、Gauss--Seidel 与 RBGS，对同一 Poisson 系统同时记录 residual norm 与已知解析解下的 error norm。
\item 人为构造低频误差 $\sin(\pi x)$ 与接近 Nyquist 的高频误差，分别做若干轮 weighted Jacobi，比较误差形状与范数。
\item 对一系列逐步加密网格，测量达到固定 residual reduction 所需的 Jacobi/GS 迭代数，并验证是否接近 $\Theta(h^{-2})$ 的趋势。
\item 扫描 $\omega\in[0.1,1.2]$，分别记录完整系统的 convergence factor 与高频 smoothing factor。说明“独立 solver 最优参数”和“multigrid smoother 最优参数”为什么可能不同。
\item 设计一个数值实验区分“这个迭代器最终收敛”和“它是一个好的 multigrid smoother”这两个命题。
\end{enumerate}
